
Empirical hypothesis testing is foundational to scientific progress, enabling researchers
to assess how well a sample represents an underlying truth.
A central challenge in this process is determining when enough data have been collected
to make a reliable decision, especially as the accuracy of statistical assessments
depends heavily on sample size.  Careful planning for statistical power and precision
is essential to avoid biased inference and ensure meaningful results \citep{cohen1988}.

Sequential hypothesis testing offers a flexible alternative to fixed-sample designs
\citep{wald1947} by allowing data to be evaluated as it accumulates.
This approach is widely used in fields such as clinical trials
\citep{jennison2000, fda2010bayesian}, quality control, A/B testing \citep{johari2022}, and financial
monitoring,
offering potential benefits in efficiency, timeliness, and ethical considerations.
Central to this approach is the distinction between two critical components.
The \textit{stopping rule} dictates when data collection ends; the
\textit{decision rule} then determines the final verdict on the hypothesis,
one of three outcomes: accept, reject, or inconclusive.
While these components are conceptually distinct, many operational methods
conflate them, triggering a stop exactly when a decision criterion is met.

This simultaneous application of stopping and decision rules is common in widely used
frequentist and Bayesian methods, such as p-value thresholds \citep{fisher1925} or posterior-based heuristics
(e.g., HDI+ROPE \citealp{kruschke2011, kruschke2013, kruschke2015doing},
Bayes Factors \citealp{kassraftery1995}).
However, if not carefully designed, this coupling can lead to confirmation bias:
premature stopping on extreme or unrepresentative early samples, a phenomenon known
as early peeking \citep{simmons2011}.
As a concrete illustration: in the fair coin setting
($\theta_{\rm true}=\theta_{\rm null}=0.5$, credible interval of 95\% and minimum effect size $\pm 0.05$),
a sequentially monitored HDI+ROPE algorithm achieves high conclusiveness
but at a systematic 6.3\% false-positive rate (Section~\ref{sec:fair_coin}).

One principled remedy is to decouple the stopping criterion entirely from the hypothesis
outcome, halting based on posterior precision rather than the direction of the result.

Following the \textit{Accuracy in Parameter Estimation} paradigm \citep{maxwell2008},
\cite{kruschke2015doing} advocated for a method that strictly decouples these two steps,
thereby addressing the risk of premature and biased conclusions from early peeking.
He proposed using a predetermined target for posterior precision as the
sole \textit{stopping rule}, independent of the hypothesis outcome. Only once this precision goal
is met is the \textit{decision rule} applied to accept or reject the null hypothesis
(via a posterior credible interval criterion).
\cite{kruschke2015doing} demonstrated that this ``Precision is the Goal'' (PitG) method
eliminates the early-peeking bias in outcome sampling. However, our analysis reveals an important
limitation: when the null hypothesis is true, or practically so, PitG often yields a high rate of inconclusive
results, rather than correctly accepting the null.\footnote{\citealt{kruschke2015doing} noted in passing that PitG ``did remain undecided
in some cases'' (\S13.3.2, page 391); our analysis shows this is a systematic and
striking effect, with the sampling budget $N_{\rm max}$ and precision goal $\omega_{\rm goal}$
as the key governing parameters. See below for quantitative examples.}

In this paper, we propose an improved variant we call
``Decisive Precision is the Goal'' (DPitG), in which precision is a necessary but no
longer sufficient stopping condition: data collection continues until \textit{both} the
precision goal and a decisive hypothesis outcome are satisfied simultaneously.
The benefit is most pronounced precisely where PitG's inconclusive rate is highest,
namely when the null is true or practically so, making DPitG particularly valuable for
confirming the null hypothesis.
While scientific incentives often favour the discovery of new effects,
rigorous confirmation of the null is equally critical in many fields: from bioequivalence
trials for generic drugs \citep{schuirmann1987, karalis2012}
to do-no-harm regression testing in software
deployment and futility stopping in programme evaluations.
Unlike \textit{Null Hypothesis Significance Testing} (NHST, see Appendix~\ref{app:nhst}),
which can only fail to reject the null, the precision-based framework enables positive
acceptance as a genuine and actionable verdict, and DPitG ensures this verdict is
reliably reached.
This fully pre-specified rule is also compatible with pre-registration standards
\citep{nosek2018} and directly addresses the design-stage imprecision implicated in
the replication crisis \citep{osc2015, gelman2014}.

We compare DPitG and PitG against the HDI+ROPE algorithm, a coupled baseline that stops as
soon as the posterior yields a decisive result, through simulations with
single-group dichotomous data.\footnote{Two other widely used approaches, NHST
(p-value-based) and Bayes Factors, are not among the three algorithms studied
here. Both lack an explicit effect-size criterion (the ROPE); NHST cannot formally
accept the null, while Bayes Factors lack a direct precision goal.
Their limitations are illustrated in Appendices~\ref{app:nhst}
and~\ref{app:bayes_factors}, respectively.}
Applied to this setting, DPitG substantially outperforms both alternatives: for
example, at $\omega_{\rm goal}=0.08$, it reduces PitG's $62.1\%$ inconclusive rate
to $2.2\%$ at a median cost of only $4.7\%$ extra samples, while preserving a zero
false-positive rate, an advantage the HDI+ROPE algorithm matches only at a $6.2\%$
false-positive cost (Section~\ref{sec:fair_coin}). These advantages hold broadly
across precision goals and true parameter values
(Sections~\ref{sec:fair_coin}--\ref{sec:trends}).
The framework extends to continuous outcomes
(Appendix~\ref{app:expected_stop}) and two-group comparisons
(Appendix~\ref{app:extensions}).
Alongside the empirical results we provide a closed-form planning formula, an online
interactive calculator, and open-source code to support adoption in applied settings.

Section~\ref{sec:methods} describes the theoretical framework and the three stopping
algorithms; Section~\ref{sec:results} presents simulation results; and
Section~\ref{sec:discussion} places the findings in methodological context.
